% SPDX-FileCopyrightText: 2016-2026 Lari Natri <lari.natri@iki.fi>
% SPDX-License-Identifier: GPL-3.0-or-later
%
% Public author controls and early initialization.
% This file is part of the 'ulsbs' package.
%
% Loaded by ulsbs.sty after package imports, with @ a letter.
% Requires etoolbox plus expl3/xparse; later styling supplies colors/fonts and
% modules supply deferred dependencies (ulSongPosition, textnote, etc.).
% Owns initial public flags, counters, font/hooks, lengths and break penalties.
% Initialization executes at package load, including chord/flat font metrics;
% definitions here are defaults for later overrides, not song/shipout handlers.
%
% Layered setup: bootstrap definitions here are followed by the shared visual
% policy in ulsbs-style.code.tex and effective defaults in the domain modules.
% ulsbs-book-defaults.code.tex supplies early book defaults and loads the
% permission/request runtime in ulsbs-visibility.code.tex at its original point.
% These newcommand hooks are not pre-load providecommand extension points:
% consumers renew them after loading the package. In particular, common style
% replaces ulLyricFontStyle, ulPrintChord and several note/index font hooks.
% Prefer ulLyricFont/ulChordFont for shared font policy; ulLyricFontStyle remains
% the distinct lyric hook. Later changes affect rendering, not the early
% font allocation or metrics measured below.
% Set allocated lengths/skips after loading with \setlength, and assign penalty
% counts directly (e.g. \ulVerseVersePenalty=200\relax). For LaTeX counters use
% \setcounter{ulSongNumber}{1}. Public boolean setters take true or false and
% obey grouping; query their real \iful... predicates rather than redefining them.
% This is an internal bootstrap module, not a configuration file to input
% directly. See WRITING_SONGBOOKS.md for the author customization contract.

\ProvidesFile{ulsbs-init.code.tex}[Public initialization]

% === Boolean controls ===

% Keep predicates as real conditionals, including in TeX's skipped branches.
% Assignment helpers stay private; the public setter follows ordinary grouping.
\newcommand*\ulsbs@newbool[4]{%
  \let#1\iffalse%
  \def#2{\let#1\iftrue}%
  \def#3{\let#1\iffalse}%
  \newcommand#4[1]{%
    \ifstrequal{##1}{true}{#2}{%
      \ifstrequal{##1}{false}{#3}{%
        \PackageError{ulsbs}{Invalid boolean value: ##1}{Use true or false.}}}}%
}

\ulsbs@newbool\ifulSlides\ulsbs@slidestrue\ulsbs@slidesfalse\ulSlides
\ulsbs@newbool\ifulTransposeCapos\ulsbs@transcapostrue\ulsbs@transcaposfalse\ulTransposeCapos
\ulsbs@newbool\ifulNoLyrics\ulsbs@nolyricstrue\ulsbs@nolyricsfalse\ulNoLyrics
\ulsbs@newbool\ifulSongIndexes\ulsbs@songindexestrue\ulsbs@songindexesfalse\ulSongIndexes\ulsbs@songindexestrue
\ulsbs@newbool\ifulSeparateIndexes\ulsbs@sepindexestrue\ulsbs@sepindexesfalse\ulSeparateIndexes\ulsbs@sepindexestrue
\ulsbs@newbool\ifulPagePreludesEnabled\ulsbs@pagepreludestrue\ulsbs@pagepreludesfalse\ulPagePreludesEnabled
\ulsbs@newbool\ifulTitleBadgeOverlap\ulsbs@titlebadgeoverlaptrue\ulsbs@titlebadgeoverlapfalse\ulTitleBadgeOverlap
% === Numbering counters and printers ===

\newcounter{ulSongNumber}
\newcounter{ulVerseNumber}
\let\ulSongNumberStyle\arabic
\let\ulVerseNumberStyle\arabic

% No arguments; expand counters through replaceable public numbering styles.
\renewcommand\theulSongNumber{\ulSongNumberStyle{ulSongNumber}}

\renewcommand\theulVerseNumber{\ulVerseNumberStyle{ulVerseNumber}}

% {label}: render song/verse labels in current layout; no counter changes.
\newcommand\ulPrintSongNumber[1]{\sffamily\bfseries\LARGE#1}

\newcommand\ulPrintVerseNumber[1]{\ulLyricFontStyle#1.\ }

% {box register}: offset and consume the already-rendered verse-number box.
\newcommand\ulPlaceVerseNumber[1]{\kern\dimexpr-\leftskip-\parindent\relax\box#1}

% === Font and rendering hooks ===

% Argument-free font declarations/hooks below affect the caller's current group.
% Empty hooks are extension points; verse justification uses shared style.
\newcommand\ulLyricFontStyle{\normalfont\normalsize}

\newcommand\ulSongTitleFont{\sffamily\ifulSlides\Huge\else\slshape\Large\fi}

\newcommand\ulVerseFont{}

\newcommand\ulEveryVerse{}

\newcommand\ulVerseJustify{\justifyleft}

\newcommand\ulSongMark{}

\newcommand\ulVerseMark{}

\newcommand\ulNoteFont{}

\newcommand\ulMeterFont{\tiny\sffamily\upshape}

\newcommand\ulEchoFont{\ifdim\fontdimen1\font>0pt\upshape\else\slshape\fi}

% {chord text}: emit styled chord ink. Empty ulChordLocals/ulChordLineParams are
% notation-local setup hooks; ulSharpSymbol emits an inline accidental glyph.
\newcommand\ulPrintChord[1]{\sffamily\slshape\large#1}

\newcommand\ulChordLocals{}

\newcommand\ulChordLineParams{}

\newcommand\ulSharpSymbol{\ensuremath{^\#}}

% No arguments; expand to color names resolved by later shared styling.
\newcommand\ulSongNumberBackgroundColor{ulSongNumberBackgroundColor}

\newcommand\ulNoteBackgroundColor{ulNoteBackgroundColor}

\newcommand\ulIndexBackgroundColor{ulIndexBackgroundColor}

% {text}: emit chord-colored musical prose under the shared lazy comment guard.
\newcommand\musicnote[1]{\ulsbs@musiccomment{\textnote{\color{ulChordColor}#1}}}

% Argument-free index font declarations, applied within display-owned groups.
\newcommand\ulIndexHeadingFont{\sffamily\bfseries\LARGE}

\newcommand\ulIndexTitleFont{\sffamily\slshape}

\newcommand\ulIndexLyricFont{\rmfamily}

\newcommand\ulIndexTagFont{\sffamily\small\slshape}

\newcommand\ulIndexAuthorFont{\small\bfseries}

\newcommand\ulIndexRefsFont{\normalfont\normalsize}

% {heading}: emit index continuation text; no collection-state effects.
\newcommand\ulIndexContinuation[1]{\small\textbf{#1} (continued)}

% === Scoped layout controls ===

% No-argument switches: suppress lyrics, request page preludes (also disable fit
% positioning), or hide numbering widths/verse printer. Dependencies resolve later.
\newcommand\nolyrics{\ulsbs@nolyricstrue}

\newcommand\ulPagePreludes{\ulsbs@pagepreludestrue\ulSongPosition{0}}

\newcommand\ulNoSongNumbers{\ulSongNumberWidth=0pt\relax}

\newcommand\ulNoVerseNumbers{\def\ulPrintVerseNumber##1{}\ulVerseNumberWidth=0pt\relax}

% No arguments; set stanza break penalties and bottom glue in the current scope.
\newcommand\ulSeparateVerses{\ulVerseVersePenalty=-10000\relax\let\ulColumnBottomGlue\@flushglue}

% No arguments; adjust paragraph metrics for note prose in its formatting group.
\newcommand\ulNoteJustify{\advance\baselineskip by 1pt\relax
  \leftskip=0pt\rightskip=0pt\parfillskip=\@flushglue\parindent=0pt}

% {original}{transposed}: default chord formatter emits only the second value.
\newcommand\ulTransposedChordFormat[2]{#2}

% No arguments; emit rights metadata at postlude construction, not at shipout.
\newcommand\ulExtendPostlude{\ulSongCopyright\ \ulSongLicense\unskip}

% === Layout registers and pre-typography numbering metrics ===

% No arguments; independent, scoped layout policy for every variant. Finite
% stretch absorbs unused column space instead of stretching compact lyric lines.
% Authors may renew it (including to zero); ulSeparateVerses selects flush glue.
\newcommand\ulColumnBottomGlue{0pt plus .5\textheight}
\newlength\ulsbs@verse@separation
\newlength\ulAfterPreludeSkip
\newlength\ulBeforePostludeSkip
\newlength\ulSongNumberWidth
\newlength\ulVerseNumberWidth
\newlength\ulIndexHeadingWidth
\newlength\ulRepeatBarWidth
\newlength\ulSongBarHeight
% Load-time numbering metrics: preserve this pre-typography font context and
% execution order, even if later styles hide verse numbers or replace printers.
\settowidth\ulSongNumberWidth{\ulPrintSongNumber{999.}}
\settowidth\ulVerseNumberWidth{\ulPrintVerseNumber{9\kern1em}}
\ulIndexHeadingWidth=1.5cm\relax
\ulRepeatBarWidth=1pt\relax\ulSongBarHeight=1pt\relax
\newcount\ulVerseVersePenalty
\newcount\ulBreakPenalty
\newcount\ulHardBreakPenalty
\newcount\ulSongPenalty
% These are document defaults, before the shared ULSBS spacing policy is applied.
\parindent=.25in\relax

% === Flat-size defaults, accidental glyphs and early chord metrics ===

\ExplSyntaxOn
\prop_new:N \g__ulsbs_song_flat_sizes_prop

% {current font size}{flat font size}: globally store expanded size strings.
% The following grouped table initialization populates defaults without leaking
% sequence scratch; entries persist for flat rendering throughout the document.
\DeclareDocumentCommand \ulDeclareFlatSize {m m}
  {\prop_gput:Nxx \g__ulsbs_song_flat_sizes_prop {#1}{#2}}
\group_begin:
\clist_map_inline:nn
  {5/5,6/6,7/6,8/7,9/8,10/9,10.95/10,12/10.95,
   14.4/12,17.28/14.4,20.74/17.28,24.88/20.74}
  {\seq_set_split:Nnn \l_tmpa_seq {/}{#1}
   \exp_args:Nxx \ulDeclareFlatSize
     {\seq_item:Nn \l_tmpa_seq {1}}{\seq_item:Nn \l_tmpa_seq {2}}}
\group_end:

% No arguments; public sharp alias, resolved through the replaceable glyph hook.
\DeclareDocumentCommand \shrp {}{\ulSharpSymbol}

% No arguments; emit a raised flat, locally selecting the size-table override
% and symbol font when available (otherwise a math flat). No font state leaks.
\DeclareDocumentCommand \ulFlatSymbol {}
  {\leavevmode\raise.5ex\hbox{\group_begin:
    \prop_get:NVNT \g__ulsbs_song_flat_sizes_prop \f@size \l_tmpa_tl
      {\fontsize{\l_tmpa_tl}{\f@baselineskip}\selectfont}
    \cs_if_exist:NTF \ulSymbolFont
      {{\ulSymbolFont\char"266D}}{\ensuremath{\flat}}\group_end:}}

% No arguments; public flat alias used by active accidentals and chord metrics.
\DeclareDocumentCommand \flt {}{\ulFlatSymbol}
% Preserve this early sample's font allocation/order for PDF range identities.
% It is not a leading strut: stanza/prose leading uses the current lyric font,
% and only retained visible annotations contribute additional ink clearance.
\setbox\@tempboxa\hbox{{\ulPrintChord{ABCDEFG\shrp\flt/j7}}}
\baselineskip=\f@size pt\relax
% Scale a dimension, not font-size text: 11pt classes use the decimal 10.95.
\advance\baselineskip 0pt~plus~\dim_eval:n {\f@size pt / 10}\relax
\ExplSyntaxOff
